% !TEX root = ../main.tex
\chapter{Money and the State: Domestic}\label{ch:state}
\begin{paracol}{2}

\begin{sources}
\begin{tabularx}{\linewidth}{@{}l l L@{}}
\toprule
\textbf{Segment} & \textbf{Length} & \textbf{Content}\\
\midrule
\seg{3}{1}{L3.1 FT: Quantitative Easing and the Fed} & 7:47 & The Fed's balance sheet 2007--2010; QE; the expected QE3.\\
\seg{3}{2}{L3.2 Allyn Young: Money and Economic Orthodoxy} & 9:09 & Where Young stands against the orthodoxy of economists and against outside opinions.\\
\seg{3}{3}{L3.3 National Banking System before the Fed} & 3:22 & The pyramid of reserves: country banks, reserve cities, New York.\\
\seg{3}{4}{L3.4 Civil War Finance, bonds and loans} & 8:49 & Bond sales, the \$150 million bank loan of 1861, suspension of gold payments.\\
\seg{3}{5}{L3.5 Civil War Finance, legal tenders} & 7:07 & The greenbacks: state money inserted into the hierarchy.\\
\seg{3}{6}{L3.6 National Banking System, origins} & 6:30 & Bonds sold to banks, bank notes backed by bonds; a fixed quantity of base money.\\
\seg{3}{7}{L3.7 National Banking System, instability} & 5:57 & Seasonal demand for currency, call loans, crises; 1907 and clearinghouse certificates.\\
\seg{3}{8}{L3.8 Federal Reserve System, plan} & 6:32 & Discounts and rediscounts: elastic reserves and elastic currency backed by loans to Main Street.\\
\seg{3}{9}{L3.9 Federal Reserve System, actual} & 6:48 & War finance fills the Fed with Treasury bills; open market operations; state money versus private money.\\
\bottomrule
\end{tabularx}

\smallskip
\textbf{Files.} Transcripts \file{materials/transcripts/L03-P0*.txt}.
\textbf{Reading.} Allyn Young's four chapters (\file{Allyn Young.pdf}, see \cref{sec:fp-young}); this lecture translates
their history of American money into balance sheets.
\end{sources}

\section*{Motivation}
Lecture~2 left the state out: the central bank was just a bankers' bank \ts{2}{2}{2:11}. Here the state comes in, through
the history that Allyn Young tells: how the United States financed the Civil War, how that war finance produced the
National Banking System, why that system was unstable, and how the Federal Reserve was designed to fix it and what it became
instead. The history serves two purposes \ts{3}{9}{3:37}:
\begin{enumerate}
  \item practice with balance sheets: every episode is a set of T-accounts;
  \item a conceptual point: the monetary system is a \emph{hybrid} of state money and private money, and a theory built on only
        one of the two misses half of it.
\end{enumerate}
A recurring pattern: \emph{the history of American central banking is a history of war finance} \ts{3}{8}{0:07}.

% ------------------------------------------------------------------
\section{FT: quantitative easing and the Fed}\label{sec:st-ft}

\begin{casestudy}[FT, September 2012: ``Fed's expected QE move weighs on the dollar'']
Global stocks rallied and the dollar was sold amid expectations that the Fed would announce a third round of quantitative
easing, and that Germany's constitutional court would allow Europe's new bailout fund (the court did, with reservations)
\ts{3}{1}{0:28}. ``QE'' is a word not yet in the textbooks: it was invented during the crisis
\ts{3}{1}{1:09}.
\end{casestudy}

\paragraph{The Fed's balance sheet in the crisis.} The chart shown in class plots the Fed's balance sheet from January 2007 to
March 2010, assets as positive numbers and liabilities as negative ones, with two vertical lines: the collapse of Bear Stearns
(March 2008) and of Lehman Brothers and AIG (September 2008) \ts{3}{1}{1:52}. It shows three phases.
\begin{enumerate}
  \item \textbf{Interest rates.} The crisis began in August 2007, and the Fed first lowered its rate from 5\% to 2\% in steps,
        trying to make the system more elastic \ts{3}{1}{2:33}. It was not enough: Bear Stearns failed \ts{3}{1}{2:53}.
  \item \textbf{Changing the composition of assets.} Before the crisis the Fed held essentially Treasury bills, against currency
        \ts{3}{1}{3:13}. After Bear Stearns it sold about half of its bills and lent the proceeds directly to banks and
        broker-dealers to support financial markets; liabilities did not change \ts{3}{1}{3:34}.
  \item \textbf{Expanding both sides.} After Lehman and AIG this was not enough either, and the Fed expanded its own liabilities:
        the balance sheet roughly doubled in a few weeks \ts{3}{1}{3:54}. First came short-term loans, to domestic
        and to international banks: the \emph{central bank liquidity swaps} with the ECB, the Bank of Japan and others reached
        about \$600 billion, because the problem had become global \ts{3}{1}{4:34}. As these facilities ran off
        around the turn of the year, the Fed started buying mortgage-backed securities, about a trillion dollars of them
        \ts{3}{1}{5:15}.
\end{enumerate}

\begin{nfigure}
\begin{tikzpicture}[font=\scriptsize,x=0.42cm,y=0.42cm]
  % stylised, not data: time on x (2007-2010), assets up, liabilities down
  \draw[-{Stealth[length=3pt]}] (0,0) -- (14.5,0) node[right]{time};
  \draw[-{Stealth[length=3pt]}] (0,-6.5) -- (0,6.5);
  \node[left] at (0,5) {assets}; \node[left] at (0,-5) {liabilities};
  % assets
  \fill[defcol!45] (0,0) -- (0,3.6) -- (5,3.6) -- (6,1.8) -- (14,1.8) -- (14,0) -- cycle;
  \fill[fincol!45] (5,3.6) -- (6,1.8) -- (8,1.8) -- (9,1.8) -- (8.5,5.5) -- (7,5.8) -- cycle;
  \fill[cscol!45] (7,5.8) -- (8.5,5.5) -- (9,4) -- (14,6) -- (14,1.8) -- (9,1.8) -- (8,1.8) -- cycle;
  % liabilities
  \fill[defcol!25] (0,0) -- (0,-3.6) -- (14,-4.2) -- (14,0) -- cycle;
  \fill[thmcol!30] (6.5,-3.7) -- (7.2,-6.2) -- (14,-6.0) -- (14,-4.2) -- cycle;
  \draw[dashed] (5,-6.5) -- (5,6.5) node[above]{Bear Stearns};
  \draw[dashed] (7,-6.5) -- (7,6.5) node[above right]{Lehman, AIG};
  \node[align=center] at (2.5,1.8) {Treasury\\bills};
  \node[align=center] at (7.2,3.3) {loans,\\swaps};
  \node[align=center] at (11.6,3.8) {MBS};
  \node[align=center] at (3,-1.8) {currency};
  \node[align=center] at (10.5,-5.1) {bank reserves};
  \node[below] at (0.5,-6.5) {2007}; \node[below] at (13.5,-6.5) {2010};
\end{tikzpicture}
\caption{Schematic of the Fed's balance sheet in the crisis, after the chart shown in class \ts{3}{1}{2:12} (shapes only
indicative, not data): first a change of composition (bills sold, loans made), then an expansion of both sides (loans, swaps,
then mortgage-backed securities, against reserves).}\label{fig:st-fedbs}
\end{nfigure}

\begin{coursenote}[QE1, QE2, QE3]
In class Mehrling calls the expansion after Lehman ``QE1'' and the purchases of mortgage-backed securities ``QE2''
\ts{3}{1}{4:14}. In the usual numbering the first program of large-scale asset purchases (QE1, from
November 2008, extended in March 2009) included the MBS purchases, about \$1.25 trillion by March 2010; QE2 (November
2010) was \$600 billion of long-term Treasuries; QE3, announced on 13 September 2012, the day after this lecture's FOMC
meeting, was open-ended MBS purchases of \$40 billion a month. Mehrling also says ``March of 07'' for Bear Stearns, which
failed in March 2008.
\end{coursenote}\begin{aside}[Why ``QE1'']
For someone ``of a certain age'', QE1 and QE2 are ocean liners, the \emph{Queen Elizabeth} and the \emph{Queen Elizabeth 2};
Mehrling guesses the joke was coined by the \emph{Financial Times}, a British paper \ts{3}{1}{4:14}.
\end{aside}

\paragraph{QE3: two ways to lower long rates.} At Jackson Hole the Fed had signalled that the economy was softening. Short
rates were already near zero; the question was how to lower the long end of the term structure further
\ts{3}{1}{5:36}:
\begin{enumerate}
  \item \textbf{Announce} that short rates will stay low for a very long time: if the market believes it, long rates fall, by the
        expectations hypothesis (\cref{def:nh-eh}) \ts{3}{1}{6:38};
  \item \textbf{buy long bonds}, expanding the balance sheet on both sides, which pushes up their prices: a capital gain for
        whoever owns them, hence the excitement of the market \ts{3}{1}{6:59}.
\end{enumerate}
In practice the Fed would probably do both \ts{3}{1}{6:59}.

\begin{definition}[New concepts: quantitative easing]
\emph{Quantitative easing} (QE) is the purchase by a central bank of large quantities of assets, especially long-term ones,
paid for by creating reserves, i.e.\ by expanding both sides of its balance sheet, when the short-term policy rate is already
near zero. ``Quantitative'' because the policy is stated as a quantity of purchases rather than as a rate.
\end{definition}

% ------------------------------------------------------------------
\section{Allyn Young against orthodoxy}\label{sec:st-young}

Young wrote in 1924, after the First World War; he had been the economist of the American delegation at Versailles, a
contemporary of Keynes but a much quieter personality \ts{3}{2}{0:50}. His encyclopedia chapters read
smoothly, and do not cite the literature they answer, so it is easy to miss that he takes position on many controversial
issues \ts{3}{2}{1:31}. Mehrling lists six.

\paragraph{Against economic orthodoxy} \ts{3}{2}{1:54}:
\begin{enumerate}
  \item \textbf{Barter.} Economics tells the story that money was invented to escape the inconvenience of barter (a trope Mehrling
        attributes to Jevons). Young: there was never a period of barter; money has been there ever since trade, it is
        ``baked into the cake'' \ts{3}{2}{1:54}.
  \item \textbf{Growth.} A developing economy needs a proper financial system. Orthodoxy then and now disagrees: in the Solow growth
        model growth comes from capital accumulation, population and productivity, technical change is a residual, and any
        growth due to finance is dropped into that residual; money is a veil to look through \ts{3}{2}{2:35}.
  \item \textbf{The currency principle.} Most economists think money gets its value from scarcity relative to demand, as in the
        neoclassical theory of value, and fear over-issue \ts{3}{2}{3:59}. Young criticises the currency principle:
        the National Banking System was built on it, limiting both note and deposit issue, and the lack of elasticity caused
        periodic crises; the Fed arose to give elasticity \ts{3}{2}{4:40}.
\end{enumerate}
Young was professor of monetary economics at Harvard and adviser to Benjamin Strong of the New York Fed, but his view was a
minority one among economists, then and now, though not inside central banks \ts{3}{2}{5:22}.

\paragraph{Against outside opinions} \ts{3}{2}{5:43}:
\begin{enumerate}[resume]
  \item \textbf{Chartalism.} Against the idea that the dollar has value because the state declares it legal tender
        \ts{3}{2}{6:04}.
  \item \textbf{Efficient markets.} Young points out how the inefficiency of the monetary system made asset prices swing apart from
        fundamental values; he is against speculation and for public policy to control bubbles \ts{3}{2}{6:45}.
  \item \textbf{Populism.} Above all he defends the Federal Reserve against those who saw it as a behemoth, and, against William
        Jennings Bryan, he wants the gold standard: the dollar is valuable because it is convertible into gold
        \ts{3}{2}{7:47}. This fits the hierarchy of lecture~2, with gold at the top. But the gold
        standard does not work perfectly by itself: it needs better financial institutions \ts{3}{2}{8:48}.
\end{enumerate}\begin{history}[Bryan's ``Cross of Gold'' (1896)]
At the Democratic national convention in Chicago, on 9 July 1896, William Jennings Bryan closed his speech for the free coinage
of silver with: ``you shall not crucify mankind upon a cross of gold''. He won the nomination but lost the election to William
McKinley; the Gold Standard Act of 1900 put the US formally on gold.\par
\emph{References:} \emph{Official Proceedings of the Democratic National Convention}, 1896; Gold Standard Act, 14 March 1900.
\end{history}\begin{history}[Young at Versailles]
In 1918--1919 Allyn Young headed the economic division of the Inquiry and then of the American Commission to Negotiate Peace
in Paris; Keynes was there for the British Treasury.\par
\emph{References:} C.~P.~Blitch, \emph{Allyn Young: The Peripatetic Economist}, Macmillan, 1995.
\end{history}

% ------------------------------------------------------------------
\section{The National Banking System before the Fed}\label{sec:st-nbs}

The issue of gold versus silver versus fiat money is left to the lectures on international money \ts{3}{3}{0:07}. Young's
chapters contain a table of the reserves of national banks before the Fed, in three groups \ts{3}{3}{0:30}:
\begin{enumerate}
  \item \textbf{Central reserve cities}: New York, Chicago, St.~Louis, required to hold 25\% of their deposits in ``lawful money''
        (currency, or gold) \ts{3}{3}{0:50}. Deposits in the table: \$825 million in New York, \$262 million in
        Chicago, \$116 million in St.~Louis; New York held \$218 million of lawful money, 26.8\% \ts{3}{3}{1:52}. To some extent
        the central bank of the US was the collection of New York banks, ``J.~P.~Morgan and his buddies'', which also faced
        Europe in international transactions \ts{3}{3}{1:11}.
  \item \textbf{Other reserve cities} (Boston, San Francisco, Kansas City, \dots): also 25\%, but partly held as deposits in New
        York, ``due from reserve agents'' \ts{3}{3}{2:12}.
  \item \textbf{Country banks}, about six thousand small banks: 15\% reserves, partly cash (\$199 million) and partly deposits in
        reserve cities or New York (\$226 million) \ts{3}{3}{2:32}.
\end{enumerate}

\begin{nfigure}
\begin{tikzpicture}[font=\small,
  b/.style={draw,rounded corners=2pt,align=center,minimum width=3.4cm,minimum height=0.8cm}]
  \node[b,fill=keycol!15] (g) at (0,3.6) {lawful money (gold, currency)};
  \node[b,fill=defcol!15] (ny) at (0,2.4) {New York banks (25\%)};
  \node[b,fill=intcol!15] (rc) at (0,1.2) {reserve city banks (25\%)};
  \node[b,fill=fincol!15] (cb) at (0,0) {country banks (15\%)};
  \draw[-{Stealth[length=4pt]}] (cb.east) to[out=0,in=0] node[right,font=\scriptsize]{deposits} (rc.east);
  \draw[-{Stealth[length=4pt]}] (rc.east) to[out=0,in=0] node[right,font=\scriptsize]{deposits} (ny.east);
  \draw[-{Stealth[length=4pt]}] (ny.east) to[out=0,in=0] node[right,font=\scriptsize]{reserves} (g.east);
  \draw[-{Stealth[length=4pt]},dashed] (cb.west) to[out=180,in=180] node[left,font=\scriptsize]{cash} (g.west);
\end{tikzpicture}
\caption{The layering of reserves in the National Banking System \ts{3}{3}{2:53}: country banks keep reserves as deposits in a
reserve city, reserve city banks keep them in New York, New York banks keep lawful money. ``Hierarchy, you can see that.''}
\label{fig:st-pyramid}
\end{nfigure}

% ------------------------------------------------------------------
\section{Civil War finance}\label{sec:st-civilwar}

The problem of war finance: the government must mobilise as many resources as it can to win the war \ts{3}{4}{0:07}.
It taxes as much as it can, it conscripts labour, and then it asks for volunteers: it sells bonds \ts{3}{4}{0:49}.
Salmon P.~Chase, Lincoln's Secretary of the Treasury, tried three methods in turn.

\subsection{Stage 1: bonds and loans}
\paragraph{Version 1: bonds sold to the public.} The government issues bonds and receives deposits, which it then spends on war
goods; within the banking system only the ownership of deposits changes, and the total quantity of deposits does not
\ts{3}{4}{1:33}.
\begin{taccts}
\tacct[2.0cm]{Government}{$+$Deposits (G)}{$+$Bonds}\quad
\tacct[2.0cm]{Private sector}{$+$Bonds\\$-$Deposits}{}\quad
\tacct[2.0cm]{Banks}{}{$-$Deposits (private)\\$+$Deposits (G)}
\end{taccts}
Spending the deposits on war goods transfers them back to the private sector \ts{3}{4}{2:56}.

\paragraph{Version 2: a loan from the banks.} If the public will not buy the bonds, or only at a high rate, the government can
borrow directly from the banking system \ts{3}{4}{3:16}. A loan is like a bond, a promise to pay at a future date,
but not a security traded in the open capital market \ts{3}{4}{3:57}. It is a swap of IOUs: the banks say to the government ``I
owe you \$150 million'' (a deposit), the government says ``I owe you \$150 million'' (the loan) \ts{3}{4}{4:18}. This
is what Chase did in August 1861 with the big New York bankers \ts{3}{4}{4:59}.
\begin{taccts}
\tacct[3cm]{Government}{$+$Deposits \$150m}{$+$Loan \$150m}\qquad
\tacct[3cm]{Banking system}{$+$Loan \$150m}{$+$Deposits (G) \$150m}
\end{taccts}
The deposits are bank money ``created from thin air'', an expansion of both sides of the balance sheet; ``you might call this
quantitative easing''. There is no central bank and no private sector in this operation \ts{3}{4}{5:20}.

\paragraph{Chase takes the gold.} Instead of spending the deposits at home, Chase withdrew them: deposits are promises to pay on
demand, and he demanded gold \ts{3}{4}{6:21}. He wanted the gold to buy war material abroad, where dollars were not
wanted; this turned out to be key, because the South had to barter cotton for European war material, while the North could pay
in gold \ts{3}{4}{7:23}.
\begin{taccts}
\tacct[3cm]{Government}{$-$Deposits\\$+$Gold}{}\qquad
\tacct[3cm]{Banking system}{$-$Gold}{$-$Deposits (G)}
\end{taccts}
With their gold reserves gone, the banks could not pay any other deposit in gold and suspended specie payments: the United States
left the gold standard \ts{3}{4}{8:03}.

\begin{history}[The loans of 1861 and the suspension]
In August 1861 the associated banks of New York, Boston and Philadelphia agreed to lend the Treasury \$150 million in three
instalments, paying in coin as the government drew on it. Chase insisted on withdrawing the coin rather than leaving it on deposit.
After the \emph{Trent} affair raised fears of war with Britain, the New York banks suspended specie payments on 30 December 1861,
followed by the Treasury.\par
\emph{References:} W.~C.~Mitchell, \emph{A History of the Greenbacks}, Chicago UP, 1903, part I ch.~2; B.~Hammond,
\emph{Sovereignty and an Empty Purse: Banks and Politics in the Civil War}, Princeton UP, 1970.
\end{history}

\begin{definition}[New concepts: specie, suspension, legal tender]
\emph{Specie} is coined metal (Latin \emph{in specie}, ``in kind''); \emph{specie payments} are the conversion of notes and deposits into
coin on demand, and their \emph{suspension} is the refusal to convert. \emph{Legal tender} is money that, by law, a creditor cannot
refuse in payment of a debt (Latin \emph{tendere}, to offer).
\end{definition}

\subsection{Stage 2: the greenbacks}
Banks would not fall for it twice \ts{3}{5}{0:08}. Chase's second method was legal tender notes, the \emph{greenbacks}. Mehrling
shows a dollar authorized by the act of 3 March 1863 (the third Greenback act), with the words ``this note is a legal tender for one
dollar'': you cannot get gold for it, only another dollar bill, and within the US you cannot refuse it in payment of a debt
\ts{3}{5}{0:29}.

The government buys war goods by paying with notes it prints \ts{3}{5}{2:56}; the private sector deposits
some of them (say 100) in banks, where they become the banks' reserves \ts{3}{5}{4:21}:
\begin{taccts}
\tacct[2.0cm]{Government}{$+$War goods}{$+$Legal tenders 400}\quad
\tacct[2.0cm]{Private sector}{$-$War goods\\$+$Legal tenders 300\\$+$Deposits 100}{}\quad
\tacct[2.0cm]{Banks}{$+$Legal tenders 100}{$+$Deposits 100}
\end{taccts}
Depositors now withdraw legal tenders, not gold. In the language of lecture~2, Chase inserted a layer of state money between gold and
deposits, at the top of the US hierarchy \ts{3}{5}{5:09}.

\paragraph{Consequences.} Legal tenders could not buy foreign goods: foreigners ask their gold value \ts{3}{5}{5:51}. Throughout the
war their value kept falling against gold, to about 50 cents on the dollar (a fall of the exchange rate), and domestic prices rose,
because the government was adding purchasing power instead of transferring it from the private sector as a bond sale does
\ts{3}{5}{6:11}. Young is ``aghast'' at this, while understanding that in war you do what it takes; the
better way is the third one \ts{3}{5}{4:01}.

\begin{coursenote}[Numbers]
In class: \$400 million authorized, ``almost 400 million'' issued, falling to ``50 cents on the dollar'' \ts{3}{5}{0:50}.
Historically the three Legal Tender Acts (25 February 1862, 11 July 1862, 3 March 1863) authorized \$450 million, of which about
\$431 million were outstanding at the peak (1864), and the greenback reached its low in July 1864, at about 35 cents in gold
(Mitchell, 1903).
\end{coursenote}\begin{aside}[Wesley Clair Mitchell]
In these passages Young draws on his close friend Wesley Clair Mitchell, author of \emph{A History of the Greenbacks} (1903),
a long-time professor at Columbia and one of the founders of the National Bureau of Economic Research (1920), the great American
institutionalist and quantitative economist \ts{3}{5}{1:55}.
\end{aside}

\subsection{Stage 3: the National Banking System}\label{sec:st-nbsorigin}
In 1863 Chase went back to bonds \ts{3}{6}{0:09}. He asked the banks to buy, at full price, bonds with a low interest
rate, paying with deposits (no gold withdrawals this time: the banks had none) \ts{3}{6}{0:58}. To sweeten the deal,
banks could \emph{issue bank notes using the bonds as collateral} \ts{3}{6}{1:39}. The government spends the deposits
on war goods; then, when the public wants to withdraw its deposits in currency, the banks print their own notes instead of paying
out their reserves \ts{3}{6}{2:20}:
\begin{taccts}
\textit{Step 1: bonds for deposits, deposits spent on war goods.}\par
\tacct[2.0cm]{Government}{$+$War goods}{$+$Bonds}\quad
\tacct[2.0cm]{Banks}{$+$Bonds}{$+$Deposits}\quad
\tacct[2.0cm]{Private sector}{$-$War goods\\$+$Deposits}{}

\medskip
\textit{Step 2: the public withdraws its deposits in bank notes.}\par
\tacct[2.0cm]{Banks}{}{$-$Deposits\\$+$Bank notes}\quad
\tacct[2.0cm]{Private sector}{$-$Deposits\\$+$Bank notes}{}
\end{taccts}
The bonds were deposited with a central authority as collateral: if the bank failed, a note holder did not need to go to the
bankrupt bank, he could get a government bond for his note \ts{3}{6}{4:48}. In effect the government borrowed from the
banks by allowing them to create money: notes one step removed from legal tenders, promises to pay legal tender, private money
\ts{3}{6}{5:30}.

\paragraph{A fixed base.} The number of such bonds was fixed, so at the end of the war they became the money supply: together with the
legal tenders and gold (after resumption), a \emph{fixed supply of base money} after 1863, which did not expand or contract with the
business cycle \ts{3}{6}{5:50}.

\begin{coursenote}[Which bonds?]
Mehrling describes ``special two percent bonds'' \ts{3}{6}{0:58}. The National Currency Act (25 February 1863) and the National Bank
Act (3 June 1864) required national banks to deposit US bonds with the Comptroller of the Currency to secure their notes (notes up to
90\% of the bonds' value); in the 1860s these were mostly the war bonds paying 5--6\%. The 2\% ``consols of 1930'', used for note
issue until 1935, date from the Gold Standard Act of 1900. Also, the 25\%/15\% reserve requirements were set in 1863--64; Chicago
and St.~Louis became central reserve cities only in 1887.
\end{coursenote}

% ------------------------------------------------------------------
\section{The instability of the National Banking System}\label{sec:st-instability}

Most deposits were in the country banks, \$2,627 million in Young's table, whose reserves were mostly deposits in the money
centres \ts{3}{7}{0:07}. The United States was an agricultural country, and the demand for currency had a strong seasonal cycle: at
harvest time farmers withdrew notes to move the crops; Young mentions \$50 million \ts{3}{7}{0:29}.

\paragraph{The seasonal squeeze} \ts{3}{7}{1:09}:
\begin{enumerate}
  \item farmers withdraw \$50 million in cash from country banks;
  \item country banks replenish their cash by drawing \$50 million on their deposits in New York;
  \item New York banks held total reserves of \$221 million: they lose 25\% of them and fall far below the required 25\%
        \ts{3}{7}{1:30};
  \item to rebuild reserves they call in their \emph{call loans} to stock-market speculators: in slack seasons they had lent idle
        reserves overnight to brokers, and at harvest they want the money back \ts{3}{7}{3:09};
  \item the speculators sell what they bought: stock prices fall, interest rates rise \ts{3}{7}{3:50};
  \item the banks also try to draw reserves from the world money market, London, disrupting the world gold market
        \ts{3}{7}{4:13}.
\end{enumerate}
``The United States is this big child in the world monetary system, and every fall it throws a little fit'' \ts{3}{7}{4:41}.

\begin{taccts}
\tacct[2.4cm]{Country banks}{$-$Deposits in NY 50}{$-$Deposits 50}\quad
\tacct[2.4cm]{New York banks}{$-$Reserves 50\\later: $-$Call loans}{$-$Deposits of country banks 50}
\end{taccts}

\begin{definition}[New concepts: call loan, clearinghouse]
A \emph{call loan} is a loan repayable on demand (``at call''), typically to a stockbroker against securities. A \emph{clearinghouse} is an
association of banks that settles the claims of its members on each other in one place, netting them (from \emph{clearing}, settling
accounts so that they are ``clear'').
\end{definition}

\paragraph{1907.} Europe eventually ``got sick of this'' and sent help, Paul Warburg, to help create a central bank; there was also
domestic agitation for one \ts{3}{7}{5:02}. The key moment was the crisis of 1907, when J.~P.~Morgan showed how to get out of a crisis:
the clearinghouse of the New York banks issued \emph{clearinghouse loan certificates}, a substitute for money, creating elastic reserves,
but privately and only for its members. The Fed's idea was to do the same for everyone, with a central bank \ts{3}{7}{5:23}.
(Clearinghouses are the subject of lecture~5.)\begin{history}[From 1907 to 1913]
The panic of October 1907 started with the failed attempt to corner United Copper and the run on the Knickerbocker Trust; J.~P.~Morgan
organised the rescue, and clearinghouses across the country issued loan certificates. The Aldrich--Vreeland Act (1908) created a National
Monetary Commission and emergency currency; Paul Warburg, a German banker who had settled in New York in 1902, became the main advocate of a
central bank. The Federal Reserve Act was signed by President Woodrow Wilson on 23 December 1913; the twelve Federal Reserve Banks opened in
November 1914.\par
\emph{References:} R.~F.~Bruner and S.~D.~Carr, \emph{The Panic of 1907}, Wiley, 2007; Federal Reserve Act, 1913.
\end{history}

% ------------------------------------------------------------------
\section{The Federal Reserve System: plan and reality}\label{sec:st-fed}

\subsection{The plan}
Attempts by the bankers themselves to create a central bank were not trusted; President Wilson got it through, just in time to use it
for First World War finance \ts{3}{8}{0:07}. Young describes three levels \ts{3}{8}{0:50}:
\begin{enumerate}
  \item \textbf{member banks}, ordinary banks making loans to farmers and merchants, ``Main Street'', by expanding their own deposits
        \ts{3}{8}{1:32};
  \item the \textbf{Federal Reserve Banks}, located in the same cities as the old reserve-city banks (Boston, San Francisco, \dots);
  \item the \textbf{Federal Reserve} itself at the top.
\end{enumerate}
What holds member banks back from lending is the fear that deposits are withdrawn and reserves run out \ts{3}{8}{2:13}. The plan removes
the fear: a bank short of reserves can get more from its Reserve Bank by using its loans as collateral (a \emph{discount loan}), or by
selling them to the Reserve Bank if they are in the right form \ts{3}{8}{2:33}.

\begin{taccts}
\tacct[2.6cm]{Member bank}{$+$Loans to Main Street\\$+$Reserves}{$+$Deposits\\$+$Discount loan}\quad
\tacct[2.6cm]{Reserve bank}{$+$Discount loan}{$+$Reserves of member bank}
\end{taccts}
``The alchemy of banking once again'': the member bank lends without fear, because its loans can be turned into reserves, and the Reserve
Bank creates the reserves from thin air by expanding its own balance sheet: \emph{elastic reserves}, unlike the National Banking System
\ts{3}{8}{3:41}. If depositors want cash instead, Reserve Banks can \emph{rediscount} at the Federal Reserve, which prints
notes: \emph{elastic currency} \ts{3}{8}{4:21}. (This is Young's description, not how it works today.)

\begin{definition}[New concepts: discount, rediscount]
To \emph{discount} a bill is to buy it before maturity for less than its face value; the difference is the interest. A central bank
\emph{discount loan} is a loan to a bank against eligible paper, and to \emph{rediscount} is to discount again paper that has already been
discounted once (by the member bank, then by the Reserve Bank). The rate charged is the \emph{discount rate}.
\end{definition}

The whole idea was a seasonal one: in an agricultural country the demand for currency and deposits fluctuates with the seasons, so the
supply of currency and of reserves must fluctuate too, backed by loans to Main Street \ts{3}{8}{5:26}. The Federal Reserve
Act says the assets of the Reserve Banks are to be these discounts; that is how Wilson sold it \ts{3}{8}{6:08}.

\subsection{The reality}
The Fed's balance sheet before the 2007 crisis held hardly any loans to Main Street: its main asset was Treasury bills \ts{3}{9}{0:08}.
This was true almost from the start, because the US went to war soon after the Fed was created, and at the end of the war the Reserve
Banks, and the member banks too, were ``stuffed full of Treasury bills'' \ts{3}{9}{0:28}. Young, writing in the 1920s,
treats this as temporary, and imagines \emph{open market operations} as a temporary way to provide elasticity \ts{3}{9}{0:48}.

\begin{definition}[New concepts: open market operations]
\emph{Open market operations} are purchases and sales of securities (classically Treasury bills) by the central bank in the open market,
paid for by creating or destroying reserves. They were developed by the New York Fed under Benjamin Strong in the 1920s \ts{3}{9}{1:34}.
\end{definition}

The importance of open market operations is a historical accident: the Fed was never supposed to lend to the government
\ts{3}{9}{2:15}. The framers had the greenbacks in mind: give the government the power to print money and it will print money; so only
loans to Main Street were to back the expansion of reserves and notes. Then came the First World War, ``the cat was out of the bag'', and
everybody got used to a Fed full of Treasury bills, so much so that in 2008 people found it weird when the Fed sold them and lent to banks,
to the private sector, in its commercial paper operations \ts{3}{9}{2:35}. (Mehrling adds that the textbook
account of open market operations changing the money supply is not how they actually work: repo is involved, lecture~7
\ts{3}{9}{1:54}.)

\subsection{State money and private money}
The conceptual point \ts{3}{9}{3:58}: the system is a \emph{hybrid}. The central bank is partly a bankers' bank; in war it is the
government's bank, helping it sell its bonds, as in both world wars, and that is its job, not a perversion of it; in normal times it
facilitates private capital markets; and it moves back and forth between these roles \ts{3}{9}{4:18}.

\begin{intuition}[Two theories, both partly true]
One can build a monetary theory in which the state prints money and everything else is a promise to pay it (chartalism), or one in which
everything is private credit and the state is just a borrower (metallism, starting from gold and private pledges). Neither is any good
alone, because the system is a hybrid that moves back and forth between the two \ts{3}{9}{4:40}. As with
discipline and elasticity (\cref{sec:nh-discipline}), Mehrling is neither chartalist nor metallist, and uses both
\ts{3}{9}{5:40}. The balance sheet approach is the bridge: debits and credits work the same way for state money and for
private money \ts{3}{9}{6:22}.
\end{intuition}

% ------------------------------------------------------------------
\flushside
\section*{Key formulas and ideas}
\begin{keyformulas}
\begin{tabularx}{\linewidth}{@{}l L@{}}
Bond sale to public & government $+$deposits/$+$bonds; only ownership of deposits changes, total deposits unchanged\\
Loan from banks & swap of IOUs: bank $+$loan/$+$deposit; both sides expand (new bank money)\\
Legal tenders & state money inserted between gold and deposits; adds purchasing power: depreciation and inflation\\
National bank notes & bank notes backed by government bonds; fixed base money, inelastic\\
Seasonal squeeze & farmers' cash demand $\to$ drain on New York reserves $\to$ call loans called $\to$ stocks fall, rates rise\\
Fed plan & discount and rediscount of loans to Main Street: elastic reserves, elastic currency\\
Fed reality & war finance: Treasury bills as main asset; open market operations\\
Hybrid & state money and private money; chartalism and metallism each partly true\\
\end{tabularx}
\end{keyformulas}

\section*{Exercises}

\begin{exercise}[Three ways to finance a war]
The government needs 100 of war goods. Write the T-accounts of government, banks and private sector when it
\begin{enumerate}[(a)]
  \item sells 100 of bonds to the public;
  \item borrows 100 from the banks and spends the deposits;
  \item prints 100 of legal tenders, of which the public deposits 40 in banks.
\end{enumerate}
In each case, what happens to the total of deposits and of currency? Which method adds purchasing power?
\end{exercise}

\begin{exercise}[The seasonal squeeze]
Using the numbers of \cref{sec:st-instability}, suppose New York banks have deposits of 825 and reserves of 221, and country banks draw 50.
\begin{enumerate}[(a)]
  \item Compute the New York reserve ratio before and after.
  \item How many call loans must be called (or gold imported) to restore 25\%?
  \item Redo (b) when a Reserve Bank discounts loans for the New York banks.
\end{enumerate}
\end{exercise}

\begin{exercise}[Bank notes backed by bonds]
A national bank buys 100 of government bonds with deposits and may issue notes up to 90\% of their value. Write its balance sheet after the
public withdraws 90 of deposits in notes. Why could the total of notes not respond to the seasonal demand for currency?
\end{exercise}

\begin{exercise}[QE in T-accounts]
The Fed buys 100 of long-term Treasury bonds from a pension fund that banks with Bank A. Write the changes in the three balance sheets and
explain why the market price of bonds is expected to rise.
\end{exercise}
